\section{Method and Experimental Design}
\label{sec:method}

\subsection{Generalizability Theory}
\label{sec:gtheory}

Generalizability theory (G-theory) treats each measurement as one observation from a universe of admissible conditions defined by the crossed facets of an evaluation design \cite{shavelson1991generalizability, brennan2001generalizability}. The \emph{object of measurement} is the entity whose consistency is of interest. In our setting, test items ($p$) serve as the object of measurement, and the goal is to determine how reliably item-level scores generalize across the methodological conditions under which they are collected.

A \emph{G-study} decomposes total observed variance into components $\sigma^2_x$ for each facet $x$ and their interactions. Each component quantifies how much variability that source contributes to the overall measurement. Large interaction components (e.g., $\sigma^2_{pM}$ for item-by-model) indicate that facets do not affect items uniformly but instead create differential difficulty profiles. We estimate variance components via Henderson Method~I, an unbiased ANOVA-based estimator for balanced designs \cite[Ch.~3]{brennan2001generalizability}, with 1,000-resample bootstrap confidence intervals \cite{efron1993bootstrap} (item-level resampling with replacement, preserving the fully-crossed structure across all other facets) \cite{li2023bootstrap}.

A \emph{D-study} projects the generalizability coefficient under hypothetical designs with different numbers of conditions per facet:
\begin{equation}
G_\text{item} = \frac{\sigma^2_p}{\sigma^2_p + \sigma^2_\text{rel}}
\label{eq:gcoeff}
\end{equation}
where $\sigma^2_\text{rel} = \sum_i \sigma^2_{p \times f_i} / n_{f_i} + \sigma^2_\text{residual} / \prod_i n_{f_i}$ aggregates all item-by-facet interaction variances, each divided by the number of conditions for the corresponding facets \cite[Ch.~4]{brennan2001generalizability}. A D-study projects reliability for a test averaging over $n_p$ items via
\begin{equation}
G(n_p) = \frac{n_p \, \sigma^2_p}{n_p \, \sigma^2_p + \sigma^2_\text{rel}}
\label{eq:dstudy}
\end{equation}
enabling independent verification of all budget projections in Table~\ref{tab:dstudy}. This prospective capability distinguishes G-theory from classical ANOVA: rather than only attributing variance post hoc, a D-study enables optimization of evaluation designs \emph{before} data collection \cite{shavelson1989gtheory}. We define \emph{item-related variance} as the sum of the item main effect ($\sigma^2_p$) and the model$\times$item interaction ($\sigma^2_{pM}$). The latter captures differential item functioning (DIF): items whose difficulty ranking changes across models \cite[Ch.~4]{brennan2001generalizability}. Because DIF is operationally an item-level property---it reflects model-specific item difficulty profiles rather than a global model ability shift---aggregating $\sigma^2_p + \sigma^2_{pM}$ provides a unified measure of item-driven variance. We report unbundled components throughout to ensure transparency.\footnote{Notation: $p$ = item, $M$ = model, $\pi$ = precision, $\delta$ = prompt, $s$ = seed, $\tau$ = temperature, $o$ = ordering, $\sigma^2_x$ = variance component for source $x$.}

\subsection{Experimental Design: MMLU}
\label{sec:exp001}

We design a fully crossed factorial experiment (exp-001) that simultaneously varies seven facets of LLM evaluation on MMLU. Table~\ref{tab:design_matrix} summarizes the complete design matrix. Figure~\ref{fig:overview} illustrates the complete pipeline from factorial design through variance decomposition to budget projection.

\paragraph{Facets.} Items ($p$): 200 MMLU questions selected via purposive stratified sampling across 10 of MMLU's 57 subjects (20 per subject, deterministic selection with arbitrarily chosen seed 42). The 10 subjects span four disciplinary domains---STEM: abstract algebra, college physics, high school mathematics, computer security; humanities: world religions, philosophy; social sciences: US foreign policy, international law; professional: professional medicine, clinical knowledge---ensuring that variance estimates are not dominated by a single knowledge domain. Resampling stability (CV = 3.1\%) was evaluated within these 10 subjects; cross-subject generalization remains untested. Models ($M$): eight open-weight 7--9B parameter instruction-tuned models spanning eight architecture families---DeepSeek-7B, Gemma-2-9B, InternLM2.5-7B, Llama-3.1-8B-Instruct, Mistral-7B-v0.3-Instruct, OLMo-2-7B, Qwen3-8B, and Yi-1.5-9B (full identifiers in Appendix~\ref{app:models})---treated as a quasi-random facet (\S\ref{sec:model_facet}). Numerical precision ($\pi$): BF16, FP16, and FP32, all evaluated with \texttt{enforce\_eager=True} for fair comparison across precision levels. Prompt template ($\delta$): six surface perturbations of the evaluation instruction \cite{sclar2024formatting}, varying option layout (lettered vs.\ numbered), instruction wording (direct vs.\ contextual), and answer prompt cue for MC benchmarks; for FF benchmarks, templates vary chain-of-thought elicitation strategy (full template text in Appendix~\ref{app:models}). Random seed ($s$): six fixed values (42, 123, 456, 789, 1024, 2048). Sampling temperature ($\tau$): three levels (0.0, 0.3, 0.7) spanning deterministic to moderately stochastic decoding. Answer ordering ($o$): four permutations of few-shot demonstration examples (for MMLU; structurally inert for other benchmarks, see \S\ref{sec:exp002}) \cite{pezeshkpour2024ordering}.

\paragraph{Sample sizes.} The cross-model BF16 balanced subset yields 460,800 records (8 models $\times$ 200 items $\times$ 6 prompts $\times$ 6 seeds $\times$ 2 temperatures $\times$ 4 orderings). A single-model analysis (Llama) spans all three precisions, yielding 259,200 binary observations across 1,296 fully balanced conditions. The BF16 subset serves as the primary cross-model sample because BF16 is the standard inference precision for these architectures.

\begin{table}[t]
\centering
\caption{Fully crossed G-theory design. Exp-001 evaluates MMLU across all facet levels; exp-002 extends to five benchmarks with reduced crossing. $^\dagger$BF16 balanced subset (2 temperatures: 0.0, 0.7) for cross-model analysis; the full single-model design uses all 3 levels.}
\label{tab:design_matrix}
\small
\resizebox{\columnwidth}{!}{%
\begin{tabular}{llrr}
\toprule
Facet & Levels & exp-001 & exp-002 \\
\midrule
Item ($p$) & Stratified sample & 200 & 200$\times$5 \\
Model ($M$) & 8 architecture families & 8 & 8 \\
Precision ($\pi$) & BF16 / FP16 / FP32 & 3 & 1 \\
Prompt ($\delta$) & Paraphrases & 6 & 6 \\
Seed ($s$) & Fixed set & 6 & 6 \\
Temperature ($\tau$) & 0.0 / 0.3 / 0.7 & 3 & 2 \\
Ordering ($o$) & Demo permutations & 4 & 4 \\
\midrule
\multicolumn{2}{l}{Total records} & 460,800$^\dagger$ & 2,304,000 \\
\bottomrule
\end{tabular}}
\end{table}


\subsection{Cross-Benchmark Extension}
\label{sec:exp002}

To test whether the variance structure observed on MMLU generalizes across evaluation formats, exp-002 extends the design to five benchmarks: MMLU, ARC-Challenge, and HellaSwag (multiple-choice), and GSM8K and MATH (free-form generation). The design crosses 8 models $\times$ BF16 $\times$ 2 temperatures (0.0, 0.7) $\times$ 6 prompts $\times$ 6 seeds $\times$ 4 orderings $\times$ 200 items per benchmark, totaling 2,304,000 records. This reduced crossing fixes precision at BF16 and uses two temperature levels to concentrate statistical power on format-associated differences. The ordering facet operates only on MMLU, where it selects among four permutations of the few-shot demonstration examples. For all other benchmarks, the prompt construction does not use the ordering parameter; ordering is structurally inert ($\sigma^2_{\text{ordering}} = 0$ by construction) and is retained solely for design balance. The two temperature levels bracket the most common evaluation configurations: deterministic ($\tau = 0.0$) and a moderate stochastic setting ($\tau = 0.7$) representative of practical deployment. A subset analysis confirms that this reduced crossing preserves the variance structure of the full design ($\Delta G = 0.0008$; Appendix~\ref{app:robustness}).


\subsection{Evaluation Protocol}
\label{sec:protocol}

All evaluations use generation-based scoring (free-text response with answer extraction) rather than log-likelihood ranking, following \citet{hua2025flaw}. Binary correctness (1/0) is the primary metric; text-exact-match is secondary. Inference uses vLLM~0.9.2 without batch-invariant kernels, preserving CUDA-level nondeterminism as a measurable variance source \cite{messina2026background, thinkingmachines2025batch}. Few-shot protocols: 5-shot MMLU/ARC, 4-shot MATH, 8-shot GSM8K, 10-shot HellaSwag.

\subsection{Binary Data Assumptions}
\label{sec:binary}

Standard G-theory assumes continuous, approximately normal outcome data \cite{shavelson1991generalizability}. Our binary correctness scores (0/1) violate this assumption. We establish robustness through simulation and estimator properties, with two additional analyses (GLMM complement, accuracy-range stability) in Appendix~\ref{app:robustness}.

Henderson Method~I produces unbiased estimates for balanced designs regardless of outcome distribution \cite[Ch.~12]{brennan2001generalizability}, and REML estimates differ by $<$0.5pp on all components. A probit simulation, GLMM complement ($\rho = 1.0$), and accuracy-range analysis all confirm robustness (Appendix~\ref{app:robustness}).

All findings in this paper are conditioned on binary correctness scoring. This measurement resolution collapses continuous generation quality into a pass/fail outcome, which may amplify item-related variance by reducing within-item response variability and mask finer-grained variance sources (e.g., partial credit, reasoning quality) that would emerge under richer scoring rubrics. The text-exact-match comparison in \S\ref{sec:single_model} ($G = 0.007$ vs.\ $0.936$) illustrates how the scoring metric can fundamentally reshape variance structure. A partial-credit analysis on GSM8K (8 models, step-level scoring) confirms this conditionality: $G_\text{item}$ drops from 0.86 to 0.77, with all eight models showing $G_\text{partial} < G_\text{binary}$ (mean $\Delta G = -0.135$; Appendix~\ref{app:robustness}), providing evidence that D-study budgets have limited transferability to finer-grained scoring rubrics.

\subsection{Model Facet Specification}

\begin{table}[t]
\centering
\caption{Models used in all experiments. RMS = RMSNorm, LN = LayerNorm.}
\label{tab:model_ids}
\scriptsize
\begin{tabular*}{\columnwidth}{@{\extracolsep{\fill}}llrl@{}}
\toprule
Short name & Family & Params & Norm \\
\midrule
Llama-3.1-8B & Llama-3 & 8B & RMS \\
Gemma-2-9B & Gemma-2 & 9B & RMS \\
Mistral-7B-v0.3 & Mistral & 7B & RMS \\
Qwen3-8B & Qwen-3 & 8B & RMS \\
Yi-1.5-9B & Yi-1.5 & 9B & RMS \\
OLMo-2-7B & OLMo-2 & 7B & RMS \\
InternLM2.5-7B & InternLM-2 & 7B & RMS \\
DeepSeek-7B & DeepSeek & 7B & RMS \\
\bottomrule
\end{tabular*}
\end{table}

Our target universe comprises open-weight, instruction-tuned models in the 7--9B parameter range (eight architecture families, Table~\ref{tab:model_ids}); results may not generalize to API-based, base, or $>$9B models. Our eight purposively selected models are treated as a quasi-random facet \cite{meredith1993measurement, putnick2016measurement}. We report $G$ under random-model assumptions, the more conservative specification: across five benchmarks, treating models as fixed yields $G_\text{fixed}$ of 0.951--0.990 (mean 0.975) versus $G_\text{random}$ of 0.863--0.896 (mean 0.882), a consistent increase of 0.077--0.113 (mean 0.093), because model$\times$item interaction (30--36\%) enters the relative error variance under random assumptions. Consequently, our D-study budgets are upper bounds---practitioners will not under-sample regardless of facet classification. The variance decomposition itself (Tables~\ref{tab:cross_model}--\ref{tab:cross_benchmark}) is invariant to this choice. Leave-one-model-out analysis preserves component rankings across all five benchmarks, and Cohen's $\kappa$ pairwise agreement (mean $\kappa = 0.33$--$0.46$) confirms genuine differential item functioning (Appendix~\ref{app:robustness}). Four-model subset consistency checks match the observed eight-model $G$ within mean $|\Delta G| = 0.035$ across all five benchmarks (\S\ref{sec:dstudy_budgets}).

For reference, $G_\text{model}$ (treating models as the object of measurement) under the recommended 3-prompt $\times$ 3-seed design with 200 items averages 0.79 across five benchmarks (range: 0.52 on ARC to 0.93 on HellaSwag). Model ranking reliability depends primarily on model$\times$prompt interaction---ARC (8.0\%) and GSM8K (8.5\%) show lower $G_\text{model}$ because prompt sensitivity reduces model discriminability. $G_\text{item}$ remains the primary metric because item reliability is the direct optimization target for evaluation design; model reliability is a downstream consequence that improves with item count.
\label{sec:model_facet}
